\documentclass[aspectratio=169]{beamer}

\usetheme{default}
\usefonttheme{serif}
\setbeamertemplate{navigation symbols}{}
\setbeamertemplate{frametitle}{%
  \vspace{0.35em}%
  {\large\bfseries\insertframetitle}\par
  \vspace{0.25em}%
  \hrule
  \vspace{0.6em}%
}

\usepackage{amsmath,amssymb,mathtools}

\title{Adaptation Effects in Two and Three Layer Networks}
\author{Akiva Lipshitz}
\date{}

\begin{document}

\begin{frame}
  \titlepage
\end{frame}

\begin{frame}{Network}
  \begin{equation*}
    f(x)
    =
    \sum_{i=1}^{N}
    a_i\,
    \operatorname{erf}(w_i\cdot x),
    \qquad
    x\in\mathbb{R}^{d}
  \end{equation*}
  \begin{equation*}
    D=\{(x_\mu,y_\mu)\}_{\mu=1}^{P}
  \end{equation*}
\end{frame}

\begin{frame}{Targets and feature integrals}
  \begin{align*}
    d\mu_x
    &=
    p(x)\,dx,
    \\[1.2ex]
    y(x)
    &=
    \mathrm{He}_1(w_*\cdot x)
    +
    \varepsilon\,
    \mathrm{He}_3(w_*\cdot x),
    \qquad
    \varphi_m(x)
    =
    \mathrm{He}_m(w_*\cdot x),
    \\[1.2ex]
    y_m
    &=
    \langle y,\varphi_m\rangle_x
    =
    \int d\mu_x\,
    y(x)\,
    \varphi_m(x),
    \\[1.2ex]
    c_m(w)
    &=
    \langle\operatorname{erf}(w\cdot x),\varphi_m\rangle_x
    =
    \int d\mu_x\,
    \operatorname{erf}(w\cdot x)\,
    \varphi_m(x).
  \end{align*}
\end{frame}

\begin{frame}{Bayesian action}
  \begin{equation*}
    p(a,W\mid D)
    \propto
    e^{-S[a,W]}
  \end{equation*}
  \begin{equation*}
    S[a,W]
    =
    \frac{1}{2\sigma_A^{2}}\lVert a\rVert^{2}
    +
    \frac{d}{2\sigma_w^{2}}\lVert W\rVert^{2}
    +
    \frac{P\chi}{2\kappa}
    \int d\mu_x\,
    \bigl(f(x)-y(x)\bigr)^{2}
  \end{equation*}
  \vspace{0.6em}
  \begin{center}
  \small
  \begin{tabular}{ll}
    \(\kappa\) & ridge \\
    \(P\) & dataset size \\
    \(N\) & width \\
    \(d\) & input dimension \\
    \(\sigma_w^{2},\ \sigma_A^{2}\) & prior variances \\
    \(\chi\) & loss scale
  \end{tabular}
  \end{center}
\end{frame}

\begin{frame}{Phenomenology of Neural Adaptation}
  \small
  \begin{equation*}
    S[w\cdot w_*]
    =
    \frac{d}{2\sigma_w^{2}}
    \left[
      (w\cdot w_*)^{2}
      -
      \frac{\sigma_w^{2}\sigma_A^{2}\chi}{d\,N}
      \Bigl(
        t(\tilde{L})\,(w\cdot w_*)
        +
        g(\tilde{L})\,(w\cdot w_*)^{3}
        +
        u(\tilde{L})\,(w\cdot w_*)^{5}
      \Bigr)^{2}
    \right]
  \end{equation*}
  \vspace{0.8em}
  with discrepancy
  \begin{equation*}
    \tilde{L}_m
    =
    \frac{P}{\kappa}\,
    \bigl\langle y-f,\ \varphi_m\bigr\rangle_{x}
    =
    \frac{P}{\kappa}
    \int d\mu_x\,
    \bigl(y(x)-f(x)\bigr)\,
    \varphi_m(x).
  \end{equation*}
\end{frame}

\begin{frame}
  \centering
  \includegraphics[width=0.96\textwidth]{%
    ../milestones/fcn2_erf_hidden_kernel/action_h0_activation_plots/RubinRegime/V_action_kappa0p07_chi.png}
\end{frame}

\begin{frame}
  \centering
  \includegraphics[width=0.9\textwidth]{%
    ../milestones/fcn2_erf_hidden_kernel/action_h0_activation_plots/RubinRegime/learnability_vs_chi_kappa0p07_P3000.png}
\end{frame}

\begin{frame}
  \centering
  \includegraphics[width=0.96\textwidth]{%
    ../milestones/fcn2_erf_hidden_kernel/action_h0_activation_plots/RubinRegime/V_action_chi1_kappa_ray.png}
\end{frame}

\begin{frame}
  \centering
  \includegraphics[width=0.9\textwidth]{%
    ../milestones/fcn2_erf_hidden_kernel/action_h0_activation_plots/RubinRegime/learnability_vs_alpha_kappa0p07_chi1.png}
\end{frame}

% \begin{frame}{Phenomenology of Adaptation in DNNs}
% \begin{align}
% V(w)
% &=
% \frac{d}{2\sigma_w^2}
% \Bigg[
% w^2
% -
% \frac{\sigma_w^2\sigma_A^2\chi}{d\,N}
% \Big(
% \bigl(t(R)+\varepsilon\,t_{\varepsilon}(R)\bigr)\,w
% \\
% &\qquad
% +
% \bigl(g(R)+\varepsilon\,g_{\varepsilon}(R)\bigr)\,w^3
% +
% \bigl(u(R)+\varepsilon\,u_{\varepsilon}(R)\bigr)\,w^5
% \Big)^2
% \Bigg].
% \end{align}
% \begin{align}
% R
% &=
% \sigma_A^2\,
% \mathbb{E}\!\left[c\,c^\top\right]
% +
% \frac{\kappa}{P}I,
% \\
% \mathbb{E}\!\left[c\,c^\top\right]
% &=
% \frac{1}{Z}
% \int
% \begin{pmatrix}
% c_1(w)^2 & c_1(w)c_3(w)\\
% c_1(w)c_3(w) & c_3(w)^2
% \end{pmatrix}
% e^{-V(w)}\,dw,
% \\
% V(w)
% &=
% \frac{d}{2\sigma_w^2}w^2
% -
% \frac{\sigma_A^2\chi}{2N}
% \bigl(y^\top R^{-1} c(w)\bigr)^2,
% \quad
% y=\begin{pmatrix}1\\ \varepsilon\end{pmatrix},
% \quad
% Z=\int e^{-V}\,dw.
% \end{align}
% \end{frame}

\begin{frame}{Kernels of the preactivations}
  Each preactivation layer carries a kernel \(K\).
  For the erf read-in,
  \begin{equation*}
    K
    =
    \sigma_A^{2}\, J,
    \qquad
    J(x,x')
    =
    \mathbb{E}\bigl[
      \operatorname{erf}(w\cdot x)\,
      \operatorname{erf}(w\cdot x')
    \bigr].
  \end{equation*}
  \begin{equation*}
    R
    =
    K+\frac{\kappa}{P}\,I,
    \qquad
    L
    =
    K\,R^{-1},
    \qquad
    f
    =
    L\,y,
    \qquad
    y-f
    =
    \frac{\kappa}{P}\,R^{-1}y.
  \end{equation*}
  Hence \(\tilde{L}=R^{-1}y\) in the \(\varphi_m\) basis.
  A three-layer net has the same object on the next preactivation.
\end{frame}

\begin{frame}{Kernel Ridge Regression in Readout Layer}
  The readout learns a feature map:
  \(\phi_i(x)=\operatorname{erf}(w_i\cdot x)\),
  with
  \begin{align*}
    \sum_{i}
    a_i\,
    \operatorname{erf}(w_i\cdot x)
    &=
    \bigl(KR^{-1} y\bigr)(x),
    \\
    &=
    \bigl(L\,y\bigr)(x).
  \end{align*}
\end{frame}

\begin{frame}{Exact fit in the eigenbasis}
  In the eigenbasis of \(K\), with \(K\varphi_m=\lambda_m\,\varphi_m\),
  \begin{equation*}
    \bigl(KR^{-1}y\bigr)(x)
    =
    \sum_m
    \frac{\lambda_m}{\lambda_m+\kappa/P}
    \,
    y_m\,
    \varphi_m(x).
  \end{equation*}
  The same basis expands the activation,
  \begin{equation*}
    \operatorname{erf}(w\cdot x)
    =
    \sum_m
    c_m(w)\,
    \varphi_m(x).
  \end{equation*}
  The readout therefore meets \(KR^{-1}y\) mode by mode,
  \begin{equation*}
    \sum_i a_i\, c_m(w_i)
    =
    \frac{\lambda_m}{\lambda_m+\kappa/P}
    \,
    y_m.
  \end{equation*}
\end{frame}

\begin{frame}{Feature learning: GFL and grokking}
\end{frame}

\begin{frame}{What the network expresses}
\end{frame}

\begin{frame}{Sample complexity}
\end{frame}

\begin{frame}{Barrier height and time}
\end{frame}

\begin{frame}{Linear adaptation}
\end{frame}

\begin{frame}{Grokking}
\end{frame}

\begin{frame}{Inner well at larger $\chi$}
\end{frame}

\begin{frame}{Cubic learnability: NNGP and rich}
\end{frame}

\end{document}
